\subsection{Further remarks on the holomorphic functional calculus}

For every subset X of C, we denote by \(X^{*}\) the image of X under complex conjugation. Let U be an open subset of C and \(g: U \to C\) a holomorphic function. The function \(f^{*}: z \mapsto \overline{g(\overline{z})}\) is then defined and holomorphic on \(U^{*}\) . The map \(f \mapsto f^{*}\) is a continuous bijection from \(\mathcal{O}(U)\) to \(\mathcal{O}(U^{*})\) such that \((f_{1}f_{2})^{*} = f_{1}^{*}f_{2}^{*}\) and \((f_{1} + f_{2})^{*} = f_{1}^{*} + f_{2}^{*}\) for \(f_{1}\) and \(f_{2}\) in \(\mathcal{O}(U)\) .

Let C be a compact subset of C. The maps \(f \mapsto f^{*}\) of \(\mathcal{O}(U)\) in \(\mathcal{O}(U^{*})\) as U ranges over the open subsets of C containing C, induce a continuous bijection of the space \(\mathcal{O}(C)\) to \(\mathcal{O}(C^{*})\) (I, p. 49, no. 1), which is also denoted \(f \mapsto f^{*}\) and which satisfies \((f_{1}f_{2})^{*} = f_{1}^{*}f_{2}^{*}\) and \((f_{1} + f_{2})^{*} = f_{1}^{*} + f_{2}^{*}\) for \(f_{1}\) and \(f_{2}\) in \(\mathcal{O}(C)\) .

PROPOSITION 3. --- Let A be a unital involutive Banach algebra. Let \(a \in \mathrm{A}\) . The spectrum of \(a^{*}\) is the image \(\mathrm{Sp}_{\mathrm{A}}(a)^{*}\) of \(\mathrm{Sp}_{\mathrm{A}}(a)\) by complex conjugation. For every \(f \in \mathcal{O}(\mathrm{Sp}_{\mathrm{A}}(a))\) , we have \(f(a)^{*} = f^{*}(a^{*})\) .

The first assertion follows from I, p. 97. By the preceding, the map \(\varphi\) of \(\mathcal{O}(\mathrm{Sp}_{\mathrm{A}}(a))\) into A defined by \(f \mapsto (f^{*}(a^{*}))^{*}\) is a continuous unital morphism of \(\mathcal{O}(\mathrm{Sp}_{\mathrm{A}}(a))\) in A such that the image of the germ in a neighborhood of \(\mathrm{Sp}_{\mathrm{A}}(a)\) of the identity function of \(\mathbf{C}\) is equal to \(a\) . Consequently, \(\varphi\) is the map \(f \mapsto f(a)\) of the holomorphic functional calculus (I, p. 74, th. 5).

Lemma 1. --- Let D be the open unit disk in C. There exists a unique holomorphic function f defined on D such that \(f(z)^{2} = 1 - z\) for all \(z \in D\) and \(f(0) = 1\) . We have \(f^{*} = f\) .

The radius of convergence of the power series

\[
\sum_ {n = 0} ^ {+ \infty} \binom {1 / 2} {n} (- z) ^ {n}
\]

is 1 and its sum \(f\) defines a holomorphic function on D (VAR, R1, p. 27, 3.2.9) taking the value 1 at 0. It satisfies \(f(x)=\sqrt{1-x}\) for all \(x \in D \cap R\) (FVR, III, p. 19), hence \(f(z)^{2} = 1 - z\) for \(z \in D\) since the difference \(f(z)^{2} - (1 - z)\) is a holomorphic function all of whose successive derivatives vanish at 0 (VAR, R1, p. 27, 3.2.5). By definition, one verifies that \(f^{*} = f\) .

Let \(g\) a holomorphic function on D such that \(g(z)^2 = 1 - z\) for all \(z \in \mathrm{D}\) and such that \(g(0) = 1\) . The function \(g\) does not vanish and the continuous function \(f / g\) on D takes values in \(\{-1, 1\}\) ; since D is connected and \(f(0) = g(0)\) , we have \(f = g\) .
% semantic-fragments: legacy-input-seam
\relax{}
